%! Author = Jonathan
%! Date = 2022-09-26

% Preamble
\documentclass[11pt]{article}
%check
% Packages
\usepackage[letterpaper, margin=1in]{geometry}


\usepackage{amsmath}
\usepackage{enumitem}
\usepackage{graphicx}
\usepackage{blindtext}
\usepackage{titlesec}
\usepackage{hyperref}


\usepackage{mathrsfs}
\usepackage{amssymb}
\usepackage{marvosym} % additional math symbols
\usepackage[table, dvipsnames]{xcolor}
\definecolor{hyperColor}{RGB}{5, 99, 183}

\hypersetup{
    colorlinks=true,
    linkcolor={hyperColor},
    citecolor={blue!50!black},
    urlcolor={hyperColor}
}

\usepackage{tocloft} % Control table of contents, figures, etc
\renewcommand{\cftsecleader}{\cftdotfill{\cftdotsep}}
\providecolors{dvipsnames*}
\usepackage{wrapfig}
\usepackage{amsfonts} %TeX fonts from the American Mathematical Society
\usepackage[usenames,dvipsnames]{color}
\usepackage{stmaryrd}
\usepackage{booktabs}
\usepackage{float}
\usepackage{floatrow}
\usepackage[clean]{svg}

\usepackage{parskip} % normal paragraph structure

\usepackage{titling}

\usepackage{biblatex}
\usepackage[english]{babel}
\usepackage[square,numbers]{natbib}
\usepackage{ulem}
\bibliographystyle{abbrvnat}


\author{Alexander, Yonatan\\
\texttt{SN: 251097927}
\and
Aziz, Mohammad Yasin\\
\texttt{SN: 251282726}
\and
Kien Tang Tran, Michael\\
\texttt{SN: 250735158}
}
\title{Therapy Support App\\
\large Group 9}

% Document
\begin{document}


    \begin{titlepage}
        \begin{center}

            \phantom{.\\}
            \vspace{1cm}

            \begin{Huge}
                \textbf{My Mental Health App}
            \end{Huge}


            \vfill

            \begin{Large}
                \textbf{Group 9}\\
            \end{Large}

            \vspace{10pt}

            \begin{table*}[h]
                \centering
                \begin{tabular}{lrr}
                    \textbf{Alexander, Yonatan}      & \textbf{SN: 251097927} \\
                    \textbf{Kien Tang Tran, Michael} & \textbf{SN: 250735158} \\
                    \textbf{Aziz, Mohammad Yasin}    & \textbf{SN: 251282726} \\
                \end{tabular}
            \end{table*}


%            \textbf{
%                Alexander, Yonatan          (SN: 251097927)   \\
%                Kien Tang Tran, Michael       (SN: 250735158)                \\
%                Aziz, Mohammad Yasin     (SN: 251282726)}

            \vfill

            \includegraphics[width=0.4\textwidth]{NewLogoEPS}   % App Logo


            \vfill


            Computer Science 4470A\\
            \textsl{The Department of Computer Science\\
            The University of Western Ontario}

            \vspace{5pt}
            \includegraphics[width=0.1\textwidth]{Stacked_CMYK}   % Western Logo
            \vspace{5pt}

            \today\\
            \textbf{Please do NOT post our group's video to the course's YouTube Channel.}

        \end{center}
    \end{titlepage}


%
%
%    \begin{enumerate}
%        \item Business requirements related to problem's root cause
%        \item Examples:
%        \begin{enumerate}
%            \item Accurate monthly sales forecast which will help us increase our sales
%            \item need to generate monthly report that indicates sales
%            \item notifications to account executives when a customer opens a problem ticket
%        \end{enumerate}
%        \item Varies based on process model
%        \item Expected:
%        \begin{enumerate}
%            \item Shippable product
%            \item some documentation
%            \item lessons learned
%            \item Why we chose one thing and not another
%            \item If we switched from one platform to another (if we switched)
%        \end{enumerate}
%        \item stakeholders:
%        \begin{enumerate}
%            \item identify asap
%            \item anyone who has any relation to the project
%            \item external - anyone who has an influence\; including people who write certain libraries and so on  (include instructor and TA as external stakeholders)
%            \item internal - programmers, maintainable, customers, users, etc [double check this point's list]
%        \end{enumerate}
%        \item success/acceptance criteria:
%        \begin{enumerate}
%            \item measurable terms of the project's outcome for the end user, customer, stakeholders, etc.
%            \item cost
%            \item timeline
%            \item Business requirements
%            \item scope
%            \item acceptance criteria: (can be only a few)
%            \begin{enumerate}
%                \item conditions the software must fulfill
%            \end{enumerate}
%        \end{enumerate}
%    \end{enumerate}


%    \maketitle
%    Cover sheet {Project title, Group members' names and SID}

    \pagebreak

    \tableofcontents

    \pagebreak


    \section{Problem Definition}\label{sec:problem-definition}
    Therapists and patients encounter many small yet persistent challenges during the therapeutic process concerning the patients' ability to self-manage between sessions.

    While this project`s scope is not large enough to address them all, the following stood out to us:

    \begin{itemize}
        \item Mood tracking between sessions.
        \begin{itemize}
            \item Usually, at the start of a session therapists ask for an update about the time since the last session.
            However, self-reporting bias usually means that if the patient is having a bad day or is in a bad mood, their view on the time since the last session is skewed, sometimes heavily.
            Further, memory can be corrupted and events can easily be forgotten.
%            \begin{itemize}
%                \item numerical (including descriptions of the numeric values, to increase clarity)
%                \item statistical (based on numerical inputs)
%                \item descriptive (``gloomy``, ``sad``, etc)
%            \end{itemize}
        \end{itemize}
        \item Medication management
        \begin{itemize}
            \item Inventory (to not run out of important medications, mood stabilizers are a good example in the context of mental health especially)
            \item Persistence (taking medications as required and without missing a day)
        \end{itemize}
        \item Symptoms
        \begin{itemize}
            \item Medical (unrelated to the patients diagnosed )
        \end{itemize}
        \item Other self care aspects, such as sleeping, eating, etc.
    \end{itemize}
%
%    \begin{enumerate}
%        \item mood tracking between sessions
%        \begin{enumerate}
%            \item numerical (including descriptions of the numeric values, to increase clarity)
%            \item statistical (based on numerical inputs)
%            \item descriptive (``gloomy``, ``sad``, etc)
%        \end{enumerate}
%        \item Medication management
%        \begin{itemize}
%            \item Inventory
%            \item Persistence (taking medications as required)
%        \end{itemize}
%        \item Symptoms
%        \begin{itemize}
%            \item Medical (unrelated to the patients diagnosed )
%        \end{itemize}
%    \end{enumerate}


    \section{Project Objective}\label{sec:project-objective}

    \subsection{General objective guidelines}\label{subsec:general-objective-guidelines}

    \begin{enumerate}[label=O\arabic*]
%        \item User can add new entries, view existing entries, delete (remove) existing entries, to:
%        \begin{enumerate}
%            \item Mood.
%            \item Medications.
%            \item Symptoms.
%        \end{enumerate}
%        \item User will be able to view the entries from the past week, at a glance, in a dedicated view.
%        \item User can update new entries to:
%        \item The system will keep track of past entries.
%        \item Patients should be able to generate a report about their mood and medication for their therapy session in one to three clicks between each therapy session.
%        Therapists should be able to see their patient's mood and medication reports in one to two clicks between each therapy session.
        \item Patients can add new entries.
        \item Patients can view existing entries.
        \item Patients can delete existing entries.
        \item User will be able to view the entries from the past week, at a glance, in a dedicated view.
        \item Learning the process of project development through application of knowledge acquired during the Computer Science 4471A Course.
        \item The creation of an intuitive and user friendly self-management mobile app (Android ). %and/or i/OS
    \end{enumerate}


    \section{Stakeholders List}\label{sec:stakeholders-list}
    \begin{itemize}
        \item Patients - Internal stakeholder
    \end{itemize}

    \subsection{Internal Stakeholders}\label{subsec:internal-stakeholders}
    \item Therapists - External stakeholder
    \item external - anyone who has an influence\; including people who write certain libraries and so on  (include instructor and TA as external stakeholders)
    \item internal - programmers, maintainable, customers, users, etc [double check this point's list]

    \begin{enumerate}
        \item Patients (main users and the target audience).
    \end{enumerate}

    \subsection{External Stakeholders}\label{subsec:external-stakeholders}

    \begin{enumerate}
        \item Therapists
        \item Legislation and public policy (for example security of medical information)
        \item Governing bodies (such as the CRPO - The College of Registered Psychotherapists of Ontario)
        \item Instructor
        \item TA
        \item Libraries:
        \begin{enumerate}
            \item Android.
            \item Google.
            \item MPAndroid (to create the graphs).
            \item JetBrains (makers of Android Studio).
            \item Java Libraries
        \end{enumerate}
    \end{enumerate}


    \section{Success/Acceptance Criteria for each Stakeholder}\label{sec:success/acceptance-criteria-for-each-stakeholder}

%    \subsection{Patients}\label{subsec:patients}
    \begin{itemize}
        \item As a user, I want to be able to store my mood on a particular day.
        \item As a user, I want to be able to rate my mood on a particular day on a scale of 0 to 10.
        \item As a user, I want to be able to store my medication details such as the medication, dosage, date, and doctor.
%        \item As a user, I want to be able to store when I take my medication.
        \item As a user, I want to be able to have an organized collection of mood and medication reports for my therapy visit.
        \item As a user, I want to be able to create a visual representation of my recent self reported mood values.
    \end{itemize}

%    \subsection{Therapist}\label{subsec:therapist}
%    \begin{itemize}
%        \item As a user, I want to be able to look at my patient's reports in an organized manner.
%        \item As a user, I want to be able to see the medication list stored on the application.
%    \end{itemize}


%    \section{Chosen Design Patterns}\label{sec:chosen-design-patterns}
%
%    \begin{enumerate}
%        \item \textbf{State} Accounts' state\cite{RefactoringGuru}
%        \item \textbf{Visitor} Output formatting\cite{RefactoringGuru}
%        \item \textbf{Strategy} graphs (possibly multiple, first one will be line plot)\cite{RefactoringGuru} %https://www.python-graph-gallery.com/
%        \item \textbf{Singleton} Interface\cite{RefactoringGuru}
%        \item \textbf{Additional creational pattern} for different `screens` (interfaces)\cite{RefactoringGuru}
%    \end{enumerate}


    \section{Use case diagram(s)}\label{sec:use-case-diagram(s)}


    \section{Selected Use case Descriptions {only two descriptions}}\label{sec:selected-use-case-descriptions}


    \section{Sequence diagram(s) {for the selected use case for descriptions}}\label{sec:sequence-diagram(s)}


    \section{System Architecture}\label{sec:system-architecture}

    Unfortunately, writing the code in Java meant that our chosen architecture (Model-View-ViewModel) was not feasible, resulting in needing to change the System Architecture to Model View Controller very late in development.

    \textbf{System Architecture: Model View Controller}.


    \section{Detailed Class diagram(s)}\label{sec:detailed-class-diagram(s)}


    \section{State-machine diagram {for the whole system, if possible}}\label{sec:state-machine-diagram}


    \section{Entity Relationship Diagram (Data modelling)}\label{sec:er--diagram-(data-modelling)}


    \section{GitHub Link}\label{sec:github-link}

    \textbf{You MUST ask Yonatan Alexander for access, as this repository is set to private.}\\

    \href{Group Therapy-Support-App}
    https://github.com/Jon-AL/4471---Therapy-Support-App


    \section{Conclusion}\label{sec:conclusion}

    Our project, in many ways, did not go according to plan.
    We made a number of mistakes, but we believe the biggest was not making the project in Kotlin.


    \section{References}\label{sec:reference}

    \begin{enumerate}
        \item \textbf{MyTherapy (app):} \qq{Your personal pill reminder and medication tracker app}\cite{MyTherapy}
        \item \textbf{Bipolar UK’s Mood Tracker (app):} \qq{Our new Mood Tracker app can make it much easier to record your daily mood, medications, emotions and how much sleep you’ve had}\cite{BiPolUK}
    \end{enumerate}

    \bibliography{main}

    \printbibliography[heading=subbibintoc]

    \appendix{Project WBS}

    \appendix{Task Assignment Matrix}

    \begin{figure}
        \label{fig:Gant Chart} %GANT CHART
    \end{figure}

    \appendix{Sample of commits on the selected version control system}


    \appendix{Wants}

    The following were not added to the system.

    \begin{enumerate}
        \item docker
        \item gant chart
        \item Reminder system
        \item Privacy
        \item Minutes
        \item Encryption for data (including for local storage)
        \begin{itemize}
            \item Was implemented as part of Android since Android 11.
        \end{itemize}
        \item Login system
        \item Therapist interface
    \end{enumerate}

    \appendix{Minutes}

    \subsection{26-Sep-’22}\label{subsec:26-sep-22}

    \textbf{Present:}
    \begin{itemize}
        \item Yonatan
        \item Michael
    \end{itemize}

    \textbf{Not Present:}

    \begin{itemize}
        \item Yasin
    \end{itemize}

    \textbf{Discussed:}

    \begin{itemize}
        \item  Development proposal and development details.
        \item  The project will be a therapy application that helps users gather, store and present their information to be used by patients.
        \item Suggested so far:
        \begin{itemize}
            \item Incremental development.
            \item Smartphone app.
            \item Web app.
        \end{itemize}
        \item  Suggested Programming languages:
        \begin{itemize}
            \item Dart.
            \item Python.
            \item Javascript.
            \item C++.
            \item Java.
        \end{itemize}
        \item Stakeholders:
        \begin{itemize}
            \item Clients.
            \item Therapists and other mental health professionals.
        \end{itemize}
        \item Features:
        \begin{itemize}
            \item Mood self-reported values.
            \item Medication tracking.
        \end{itemize}

    \end{itemize}

    \subsection{30-Sep-’22}\label{subsec:30-sep-22}

    \textbf{Present:}
    \begin{itemize}
        \item Yonatan
        \item Michael
    \end{itemize}

    \textbf{Not Present:}

    \begin{itemize}
        \item Yasin
    \end{itemize}

    \textbf{Discussed:}

    \begin{itemize}
        \item Project proposal was accepted.
        \item Instructor mentioned a number of points during lecture in the context of our project:
        \begin{itemize}
            \item Patient data requirements with respect to privacy.
            \item Identifying the stakeholders with who should be able to access data.
            \item We need to identify the use cases carefully of what we are able to do.
            \item We have to be able to clearly define the project objective and description.
        \end{itemize}
        \item An aggressive timeline has been suggested to keep with coursework.
        This timeline is subject to change due to other courses requiring our attention.
        \item Github was setup by Yonatan.
        \item Report is in progress.
    \end{itemize}

    \subsection{04-Oct-’22 and 05-Oct-’22}

    \textbf{Discussed:}

    A timeline has been suggested to be finishing up to the user acceptance criteria.
    QT has been suggested as a platform to build an Android app.

    The report has been started in LaTex by Yonatan.
    Proposed statements will be made by the group once we reviewed the work.

    \subsection{14-Oct-’22}

    \textbf{Present:}
    \begin{itemize}
        \item Yonatan
        \item Michael
        \item Yasin
    \end{itemize}

    \textbf{Discussed:}

    Debate about the code is that we have no firm agreements on project libraries and frameworks.
    Suggestions include: salesforce as the backend.
    Flutter or ReactNative has also been suggested.
    Qt.

    \subsection{21-Oct-’22 to 28-Oct-’22}\label{subsec:21-oct-22-to-28-oct-22}


    A lot of project will be developed on the reading week.
    Professor wants us to use Docker, and create a Gantt chart.
    He explicitly states that he wants us to assign backup roles to our task assignment matrix.

    \subsection{01-Nov-’22}\label{subsec:01-nov-22}

    \textbf{Present:}
    \begin{itemize}
        \item Yonatan
        \item Michael
    \end{itemize}

    \textbf{Not Present:}

    \begin{itemize}
        \item Yasin
    \end{itemize}

    \textbf{Discussed:}

    \item    Planned architecture:
    \begin{itemize}
        \begin{itemize}
            \item repository style for the databases.
            \item Event-driven for reminders.
            \item Component-based.
        \end{itemize}
        \item Today's objective:
        \begin{itemize}
            \item Architecture and design patterns were discussed.
            \item Get the backend done behind the scenes.
            \item Another meeting: 9:30 am to 11:30 am tomorrow (November 2nd, 2022), want to start writing code.
        \end{itemize}
    \end{itemize}

    \subsection{02-Nov-’22}\label{subsec:02-nov-222}

    \textbf{Present:}
    \begin{itemize}
        \item Yonatan
        \item Michael
        \item Yasin
    \end{itemize}

    \textbf{Discussed:}
    \begin{itemize}
        \item We are switching the development process from incremental to a hybrid agile-iterative method.
        \item Scope and requirements changed significantly throughout the process.
        \item We are also incorporating agile methods to respond to changes over time.
        \item Kanban board has been created inside GitHub repository.
        \item User stories will be created for UI to make the interface clearer to write components and understand the problems.
        \item Drafts have been started but significant changes need to be created.
        \item An Android application will be created.
        \item The data will be stored locally with some encryption to get us started.
        \item Yonatan wants to push simple skeleton code, will require approval from Michael and Yassine.
        \item We want to finish baseline features executed and finished by November 6th.
        \item The list of things we want executed:
        \begin{itemize}
            \item A basic UI - covered by Yassine
            \item Basic backend - covered by Yonatan
            \item Database integration - covered by Michael
        \end{itemize}
        \item Significant research must be implemented for any implementation to occur.
        \item No meeting will be scheduled for tomorrow.
        \item A few mini-meetings between Friday to Sunday.
    \end{itemize}

    \subsection{02-Nov-’22}\label{subsec:02-nov-22}

    \textbf{Present:}
    \begin{itemize}
        \item Yonatan
        \item Michael
    \end{itemize}

    \textbf{Not Present:}

    \begin{itemize}
        \item Yasin
    \end{itemize}

    \textbf{Discussed:}
    \begin{itemize}
        \item Backend talk
        \item Significant modifications need to be executed for medication.
        \item We will have to split this into 3 main Tables:
        \begin{enumerate}
            \item Medication info
            \item Medication inventory history
            \item Medication intake history
        \end{enumerate}
        \item Mood will be generated next.
        \item Symptoms will be generated last.
        \item Backend sample is expected to be done by Sunday late afternoon.
        \item Yassine will present a frontend demo on November 6th.
    \end{itemize}

    \subsection{14-Nov-’22}\label{subsec:14-nov-22}

    \textbf{Present:}
    \begin{itemize}
        \item Yonatan
        \item Michael
        \item Yasin
    \end{itemize}

    \textbf{Discussed:}
    \begin{itemize}
        \item  Front end and back end are being implemented.
        We are integrating today.
        \item  Our approach has changed.
        We have recognized that the Model-View-ViewModel (MVVM) pattern was the best representation.
        \item Things to do:
        \begin{itemize}
            \item Integrate front-end and back-end.
            \item Revise databases.
            \item Check on scope and extra features.
            \item Testing and presentation notes to build.
        \end{itemize}
    \end{itemize}

    \subsection{24-Nov-’22}\label{subsec:24-nov-22}
    \textbf{Present:}
    \begin{itemize}
        \item Yonatan
        \item Michael
    \end{itemize}

    \textbf{Not Present:}

    \begin{itemize}
        \item Yasin
    \end{itemize}

    \textbf{Discussed:}

    \begin{itemize}
        \item
        \item Presentation is being worked on.
        \item Yonatan changed color scheme and logo.
        \item Things to do:
        \begin{enumerate}
            \item  He wants a better option for moods.
            A better option would be spinner or radio button.
            A perfect solution would be adding a slider [Seek bar].
            \item  We need to focus more on accessibility.
            \item  Add statistics.
        \end{enumerate}
    \end{itemize}

    \subsection{27-Nov-’22}\label{subsec:27-nov-22}
    \textbf{Present:}
    \begin{itemize}
        \item Yonatan
        \item Michael
    \end{itemize}

    \textbf{Not Present:}

    \begin{itemize}
        \item Yasin
    \end{itemize}

    \textbf{Discussed:}
    \begin{itemize}
        \item Code must refactored to meet architectural standards.
        \item Code will be redone to the exact standard to the specified report
        \item Mood will be the focus.
        \item Medication and Symptoms are not the primary objective.
        \item Deadline to finish the code is tomorrow night.
        \item Deadline should be finished by Wednesday night.
        \item Deadline to finish the document components is Thursday.
    \end{itemize}

    \subsection{28-Nov-’22}\label{subsec:28-nov-22}

    \textbf{Present:}
    \begin{itemize}
        \item Yonatan
        \item Michael
    \end{itemize}

    \textbf{Not Present:}

    \begin{itemize}
        \item Yasin
    \end{itemize}

    \textbf{Discussed:}
    \begin{itemize}
        \item Multiple meetings today.
        \item We have to go with MVC. MVVM was not feasible given the circumstances.
        \item Code is being redone to make the sections better.
        \item We want to get the code ready for deployment soon.
        \item Things to do:
        \begin{itemize}
            \item Clean up redundant code
            \item add comments
            \item fix anything that seems broken.
            \item Deploy on Docker.
            \item Another addition of report documents and diagrams.
        \end{itemize}
    \end{itemize}


\end{document}
